\documentclass[11pt]{article}

\usepackage[utf8]{inputenc}
\usepackage[T1]{fontenc}
\usepackage{lmodern}
\usepackage{geometry}
\usepackage{amsmath,amssymb,amsfonts,amsthm,mathtools}
\usepackage{bm}
\usepackage{enumitem}
\usepackage{microtype}
\usepackage{hyperref}
\usepackage{titlesec}
\usepackage{booktabs}
\usepackage{array}
\usepackage{setspace}
\usepackage{mathrsfs}
\usepackage{physics}

\geometry{margin=1in}
\hypersetup{colorlinks=true, linkcolor=black, urlcolor=blue, citecolor=black}
\setlist[enumerate]{topsep=0.4em,itemsep=0.35em}
\setlist[itemize]{topsep=0.3em,itemsep=0.25em}
\titleformat{\section}{\large\bfseries}{\thesection.}{0.5em}{}
\titleformat{\subsection}{\normalsize\bfseries}{\thesubsection.}{0.5em}{}
\setstretch{1.06}

\newcommand{\Aether}{\AE{}ther}
\newcommand{\Aflow}{\AE{}ther-flow}
\newcommand{\TheoryName}{\Aflow{} Interpretation of Relativity}
\newcommand{\TheoryProgram}{\Aether{} / \Aflow{} framework}
\newcommand{\Sub}{\mathcal{A}}
\newcommand{\BridgeMetric}{\mathfrak{g}}

\newtheorem{definition}{Definition}
\newtheorem{proposition}{Proposition}
\newtheorem{remark}{Remark}
\newtheorem{corollary}{Corollary}

\title{\textbf{The \TheoryName{}}\\[0.4em]
\large First Nonminimal Candidate Branch: Derivative-Mixing Completion of the \(S/Q\) Sector}
\author{}
\date{}

\begin{document}

\maketitle

\begin{abstract}
The minimal stable-sign candidate action of the derivational bridge has now been shown to possess no nondegenerate flat-background exact-closure branch. The next task is therefore not to restate the benchmark theory, but to define the first nonminimal candidate branch worth testing.

This note performs that branch-selection step in a disciplined way. It first states the admissibility criteria any new candidate branch must satisfy if it is to remain answerable to the exact-closure benchmark of \TheoryName{}. It then compares the smallest nonminimal completion families available after the no-go result and selects the derivative-mixing completion of the \(S/Q\) sector as the first branch to test. In that branch one adds the operator
\[
\Delta\mathcal L_{\rho}
=
-\rho_I P^{AB}\partial_A S\,D_BQ^I
\]
to the minimal auxiliary sector while leaving the bridge metric, the Einsteinian benchmark, and universal matter coupling unchanged.

At quadratic order this completion changes the low-momentum scalar cancellation problem from the minimal positive-sum form to a generalized operator
\[
\mathcal B_{\rho}(k^2)
=
\kappa_Sk^2+\mu_S^2
-\bigl(\mu_{SI}+\rho_Ik^2\bigr)
\bigl(k^2\delta+\widehat M\bigr)^{-1\,IJ}
\bigl(\mu_{SJ}+\rho_Jk^2\bigr),
\]
so that
\[
\mathcal B_{\rho,2}
=
\kappa_S
+\mu_{SI}\widehat M^{-2\,IJ}\mu_{SJ}
-2\mu_{SI}\widehat M^{-1\,IJ}\rho_J.
\]
The minimal no-go theorem therefore no longer applies in its original form. This note does not prove that the derivative-mixing branch succeeds. It defines the first nonminimal branch cleanly enough that the next quadratic consistency pass can test it directly.
\end{abstract}

\section{Introduction}

The active manuscript sequence of \TheoryName{} already contains the exact-closure theory statement, its effective dynamics, its consistency analysis, its relativistic recovery statement, its flow-geometry interpretation, the substrate-kinematics bridge manuscript, the coefficient-level linearized consistency pass, and the minimal stable-sign no-go result. The scientific situation is therefore sharper than before. The minimal stable-sign action is no longer an unresolved maybe. It is obstructed as a nondegenerate exact-closure branch.

That shifts the derivational program into a new stage. If the program is to continue beyond the exact-closure benchmark, it must now nominate a concrete nonminimal branch and make clear why that branch is worth testing before any further calculations are undertaken. The present note performs exactly that task.

\paragraph{Framework and claim boundary.}
\TheoryProgram{} denotes the broader \Aether{} / \Aflow{} framework built on the claims that the \Aether{} is the underlying four-dimensional substrate of reality, the \Aflow{} is its intrinsic ordered motion, observed three-dimensional space is the local experiential slice of that deeper substrate, S-time is the experienced order of change arising from matter, light, and the \Aflow{}, and observed expansion is the three-dimensional appearance of deeper four-dimensional ordered motion. Within that framework, \TheoryName{} denotes the exact-closure theory in which the effective gravitational dynamics are adopted to be exactly Einsteinian with universal matter coupling. In the active sequence, \emph{adoption} means use of that established relativistic dynamics without claiming substrate derivation, while \emph{derivation} is reserved for a first-principles recovery from explicit substrate variables.

The scope is deliberately narrow. This note does not claim that the new branch already works. It does not claim a completed Einsteinian derivation. It defines the first nonminimal candidate branch cleanly enough that the next quadratic and curved-background tests can be posed without ambiguity.

\section{Admissibility Criteria for a New Candidate Branch}

\begin{definition}[Admissible post-no-go candidate branch]
A post-no-go candidate branch is admissible for the active derivational program only if it satisfies all of the following.
\begin{enumerate}
    \item \textbf{Exact-closure answerability.} The branch must retain \TheoryName{} as the exact benchmark and must possess a well-defined limit in which the observer-accessible effective theory contracts back to the unique Einsteinian metric sector.
    \item \textbf{Single-metric matter coupling.} Matter and light must continue to couple through the bridge metric alone, with no unsuppressed direct low-energy coupling to extra substrate variables.
    \item \textbf{Health discipline.} The branch must be testable under sign, ghost, gradient, and mode-control conditions at least as strong as those imposed on the minimal stable-sign branch.
    \item \textbf{Minimal conceptual drift.} The new structure should change only the part of the auxiliary sector needed to evade the established obstruction and should not rewrite the public theory claim.
    \item \textbf{Computable next test.} The modified branch must admit an immediate flat-background quadratic consistency pass rather than only a speculative verbal motivation.
\end{enumerate}
\end{definition}

\begin{remark}
These criteria are methodological, not yet success conditions. They determine whether a proposed branch deserves to enter the manuscript line as a real candidate worth testing.
\end{remark}

\section{Smallest Nonminimal Families Worth Considering}

After the minimal stable-sign no-go result, the smallest nonminimal completion families worth considering are the following.
\begin{enumerate}
    \item \textbf{Derivative mixing in the \(S/Q\) sector.} Add operators that mix the spatial derivatives of \(S\) and \(Q^I\) while leaving the bridge metric and matter route unchanged.
    \item \textbf{Constraint-assisted \(U\)-sector completion.} Add new auxiliary structure in the longitudinal \(U\)-sector so that the scalar cancellation problem is altered through constraint dynamics rather than through direct sign changes.
    \item \textbf{Higher-derivative completion.} Add operators that postpone the leading scalar obstruction to higher derivative order while keeping the low-energy limit controlled.
\end{enumerate}

The first family is the correct initial target for three reasons. First, it is the smallest local deformation of the auxiliary sector that directly changes the low-momentum scalar coefficient responsible for the no-go result. Second, it does so without modifying the bridge metric, the universal matter-coupling route, or the public exact-closure benchmark. Third, it remains immediately computable at the same quadratic level as the existing linearized-consistency manuscript.

\begin{proposition}[First branch-selection result]
The first nonminimal candidate branch worth testing after the minimal stable-sign no-go result is the derivative-mixing completion of the \(S/Q\) sector.
\end{proposition}

\begin{proof}
Among the three families above, only the derivative-mixing family alters the established scalar obstruction at leading low-momentum order while keeping the bridge metric, the matter route, and the benchmark claim unchanged and still allowing an immediate coefficient-level quadratic calculation. It therefore best satisfies the admissibility criteria with the smallest conceptual drift.
\end{proof}

\section{Derivative-Mixing Branch Definition}

The minimal candidate action is retained, but the auxiliary sector is extended by
\begin{equation}
\Delta\mathcal L_{\rho}
=
-\rho_I P^{AB}\partial_A S\,D_BQ^I,
\label{eq:Lrho}
\end{equation}
where \(\rho_I\) is a constant coefficient vector on the chosen background branch.

The resulting candidate auxiliary sector is
\begin{equation}
\mathcal L_{\mathrm{aux}}^{(\rho)}
=
\mathcal L_{\mathrm{aux}}
+\Delta\mathcal L_{\rho}.
\label{eq:Lauxrho}
\end{equation}

\begin{definition}[Derivative-mixing branch]
The derivative-mixing branch is the branch of the candidate derivational program defined by the action obtained from the minimal bridge action after replacing \(\mathcal L_{\mathrm{aux}}\) with \(\mathcal L_{\mathrm{aux}}^{(\rho)}\).
\end{definition}

\begin{remark}
The bridge metric \(\BridgeMetric_{AB}=G_{AB}-2U_AU_B\), the Einsteinian benchmark, and the universal matter-coupling route through \(S_{\mathrm{matter}}[\Xi^*\BridgeMetric,\psi]\) remain unchanged on this branch. Only the auxiliary \(S/Q\) sector is modified.
\end{remark}

\subsection{Exact-Closure Reversion Limit}

The intended reversion limit is
\begin{equation}
\rho_I \to 0.
\label{eq:rholimit}
\end{equation}
In that limit the derivative-mixing branch contracts exactly back to the minimal candidate action. The exact-closure benchmark of \TheoryName{} is therefore preserved as the branch-selection floor rather than rewritten.

\section{Quadratic Scalar Sector of the New Branch}

On the same admissible homogeneous background used in the minimal quadratic analysis, \eqref{eq:Lrho} contributes at quadratic order
\begin{equation}
\Delta\mathcal L_{\rho}^{(2)}
=
-\rho_I\,\partial_i\sigma\,\partial_i q^I.
\label{eq:Lrho2}
\end{equation}

The scalar \(Q\)-sector equation is correspondingly modified from the minimal form to
\begin{equation}
\mathcal A_{IJ}(-\nabla^2)\,q^J
=
-\bigl(\mu_{SI}+\rho_I(-\nabla^2)\bigr)\sigma,
\label{eq:qeqrho}
\end{equation}
where
\begin{equation}
\mathcal A_{IJ}(-\nabla^2)
\equiv
(-\nabla^2)\delta_{IJ}+\widehat M_{IJ}.
\label{eq:Aoperator}
\end{equation}
Thus
\begin{equation}
q^I
=
-\mathcal A^{-1\,IJ}(-\nabla^2)\,
\bigl(\mu_{SJ}+\rho_J(-\nabla^2)\bigr)\sigma.
\label{eq:qsolverho}
\end{equation}

Substituting \eqref{eq:qsolverho} back into the scalar sector yields
\begin{equation}
\mathcal L_{\sigma,\mathrm{eff}}^{(2)}
=
-\frac{m_S^2}{2}\bigl(\partial_0\sigma-\sigma_0\phi\bigr)^2
-\frac{1}{2}\sigma\,\mathcal B_{\rho}(-\nabla^2)\,\sigma,
\label{eq:Lsigmaeffrho}
\end{equation}
with generalized scalar operator
\begin{equation}
\mathcal B_{\rho}(-\nabla^2)
=
\kappa_S(-\nabla^2)
+\mu_S^2
-\bigl(\mu_{SI}+\rho_I(-\nabla^2)\bigr)
\mathcal A^{-1\,IJ}(-\nabla^2)
\bigl(\mu_{SJ}+\rho_J(-\nabla^2)\bigr).
\label{eq:Brhooperator}
\end{equation}

In Fourier space,
\begin{equation}
\mathcal B_{\rho}(k^2)
=
\kappa_Sk^2+\mu_S^2
-\bigl(\mu_{SI}+\rho_Ik^2\bigr)
\bigl(k^2\delta+\widehat M\bigr)^{-1\,IJ}
\bigl(\mu_{SJ}+\rho_Jk^2\bigr).
\label{eq:Brhok}
\end{equation}

Expanding at small \(k\),
\begin{equation}
\mathcal B_{\rho}(k^2)
=
\mathcal B_{\rho,0}
+\mathcal B_{\rho,2}k^2
+\mathcal O(k^4),
\label{eq:Brhoexpand}
\end{equation}
with
\begin{align}
\mathcal B_{\rho,0}
&=
\mu_S^2-\mu_{SI}\widehat M^{-1\,IJ}\mu_{SJ},
\label{eq:Brho0}
\\
\mathcal B_{\rho,2}
&=
\kappa_S
+\mu_{SI}\widehat M^{-2\,IJ}\mu_{SJ}
-2\mu_{SI}\widehat M^{-1\,IJ}\rho_J.
\label{eq:Brho2}
\end{align}

\begin{proposition}[Why the minimal no-go theorem no longer applies directly]
On the derivative-mixing branch, the low-momentum coefficient \(\mathcal B_{\rho,2}\) is no longer the sum of manifestly nonnegative terms. The new cross term
\begin{equation}
-2\mu_{SI}\widehat M^{-1\,IJ}\rho_J
\label{eq:cross}
\end{equation}
can compete with the positive minimal contributions without requiring \(\kappa_S<0\) or an indefinite \(\widehat M_{IJ}\).
\end{proposition}

\begin{proof}
Equation \eqref{eq:Brho2} differs from the minimal coefficient exactly by the cross term \eqref{eq:cross}. That term has no fixed positive sign. Therefore the minimal positive-sum argument used in the stable-sign no-go theorem does not transfer unchanged to the derivative-mixing branch.
\end{proof}

\begin{remark}
This is not yet a viability proof. It is only the reason the derivative-mixing branch deserves to replace the minimal action as the next branch to test.
\end{remark}

\section{Health Conditions That Must Still Be Passed}

The derivative-mixing branch is not automatically admissible merely because it escapes the old algebraic obstruction. It must still pass its own health gates.

The first local gate is the spatial-gradient block of the \((\sigma,q^I)\) sector. At quadratic order that block is
\begin{equation}
\mathbb K_{\mathrm{grad}}
=
\begin{pmatrix}
\kappa_S & \rho_J \\
\rho_I & \delta_{IJ}
\end{pmatrix}.
\label{eq:Kgrad}
\end{equation}
A basic stable-sign requirement is therefore
\begin{equation}
\mathbb K_{\mathrm{grad}} \succeq 0.
\label{eq:Kpositive}
\end{equation}
Equivalently, because the \(Q\)-gradient block is the identity in the present branch choice,
\begin{equation}
\kappa_S-\rho_I\delta^{IJ}\rho_J \ge 0.
\label{eq:Schur}
\end{equation}

\begin{remark}
Condition \eqref{eq:Schur} keeps the derivative-mixing branch inside the same sign-discipline philosophy as the minimal stable-sign branch while still permitting nonzero \(\rho_I\).
\end{remark}

Further tests remain mandatory:
\begin{enumerate}
    \item ghost and kinetic-sign control in the full quadratic sector
    \item absence of unsuppressed extra observer-accessible poles
    \item preservation of single-metric matter coupling
    \item compatibility with local relativistic recovery and the exact-closure benchmark
    \item curved-background continuation beyond the Minkowski patch
\end{enumerate}

\section{Immediate Research Tasks for This Branch}

Now that the branch has been chosen and written explicitly, the next calculations are fixed.
\begin{enumerate}
    \item derive the full quadratic auxiliary Lagrangian of the derivative-mixing branch, including any corrections to the longitudinal \(U\)-sector elimination
    \item determine the exact cancellation conditions replacing the minimal \(\mathcal B_0=\mathcal B_2=0\) test
    \item check whether a nondegenerate flat-background exact-closure branch exists subject to \eqref{eq:Kpositive}
    \item test whether the branch retains a healthy spectrum and controlled matter coupling
    \item only after that, extend the analysis to slowly varying and then global admissible backgrounds
\end{enumerate}

\section{Conclusion}

The minimal stable-sign action no longer serves as a serious open candidate branch. After the no-go result, the derivational program needed a new branch worth testing, not a repetition of the benchmark theory. This note provides that next step.

The first branch worth testing is the derivative-mixing completion of the \(S/Q\) sector. It is the smallest nonminimal deformation that changes the leading scalar cancellation problem while preserving the bridge metric, the matter route, and the exact-closure benchmark of \TheoryName{}. Its generalized coefficient \(\mathcal B_{\rho,2}\) contains a new cross term that removes the old minimal positive-sum obstruction.

That is enough to promote the branch into the next calculation stage, but not enough to declare success. The derivative-mixing branch now has to survive the same discipline demanded elsewhere in the active sequence: explicit quadratic coefficient control, health of the observer-accessible spectrum, and eventual curved-background continuation. Until it does, exact closure remains the scientific benchmark and reversion theory of the project.

\end{document}
